% SOURCE_STATEMENT: "After finite-time formation of a Dirac mass at $\\omega=0$ in the Nordheim equation, determine whether the solution can be continued as $f(t,\\omega)=n(t)\\delta_0/\\sqrt\\omega+h(t,\\omega)$ with $h(t,\\cdot)$ a function."
% Preparation and final self-audit: CONTRIBUTING.md prompts C and D; source/audit artifacts retained in the local physics_statements_2026_10_11 batch.
% !TeX program = lualatex
\documentclass[11pt]{article}
\usepackage[margin=25mm]{geometry}
\usepackage{fontspec}
\setmainfont{Latin Modern Roman}
\usepackage{amsmath,amssymb,mathtools}
\usepackage{unicode-math}
\setmathfont{Latin Modern Math}
\usepackage{xurl,hyperref}
\hypersetup{colorlinks=true,urlcolor=blue}
\setcounter{secnumdepth}{0}
\setlength{\parindent}{0pt}
\setlength{\parskip}{5pt}
\setlength{\emergencystretch}{3em}
\newcommand{\sref}[2]{\hyperlink{source:#1:#2}{[S#2]}}
\newcommand{\cataloguescope}[1]{\paragraph{Recorded target.}#1}
\newcommand{\R}{\mathbb R}
\newcommand{\C}{\mathbb C}
\newcommand{\N}{\mathbb N}
\newcommand{\Z}{\mathbb Z}
\newcommand{\Q}{\mathbb Q}
\newcommand{\F}{\mathbb F}
\newcommand{\eps}{\varepsilon}
\begin{document}
% UnsolvedMath ID 30005626; source catalogue code OWR-14297737-003
\setcounter{section}{30005625}
\section{Continuation of Nordheim Condensation Beyond Dirac-Mass Formation}\label{30005626}
{\small UnsolvedMath ID 30005626 · Mathematical Physics}
\subsection{Definitions and mathematical statement}

\paragraph{Shared notation.}$\mathbb N=\{1,2,\ldots\}$; the standard number fields have their usual meanings. All additional conventions are stated locally.


The Nordheim equation describes a spatially homogeneous Bose gas through an energy distribution. The energy variable is $\omega\geq0$. Condensation corresponds to a measure with an atom at zero energy; $n(t)$ is its mass. In the source normalization the particle measure uses the weight $\sqrt\omega\,d\omega$, so the expression $\delta_0/\sqrt\omega$ is a symbolic way to describe that atom, not ordinary multiplication of distributions.

\paragraph{Catalogue wording retained for traceability.}
After finite-time formation of a Dirac mass at $\omega=0$ in the Nordheim equation, determine whether the solution can be continued as $f(t,\omega)=n(t)\delta_0/\sqrt\omega+h(t,\omega)$ with $h(t,\cdot)$ a function.

\paragraph{Scope and source qualification.}A measure-valued weak formulation is necessary at the atom. Conservation laws, positivity and any uniqueness claim must be justified in that formulation.
\subsection{Short English statement}
Determine whether the Nordheim evolution can continue after condensation as a zero-energy atom plus a regular energy-density component.
\subsection{Sources}
\begin{itemize}
\item[\hypertarget{source:30005626:1}{S1}] ULAM.AI and the UnsolvedMath Contributors, UnsolvedMath, record 30005626 (OWR-14297737-003). Catalogue curation under CC BY 4.0; definitions and English exposition added locally.
\url{https://huggingface.co/datasets/ulamai/UnsolvedMath}.
\item[\hypertarget{source:30005626:2}{S2}] Classical and Quantum Mechanical Models of Many-Particle Systems (2023), Oberwolfach Reports, DOI 10.4171/owr/2023/38, pp. 2179--2246. Original report and mathematical context.
\url{https://doi.org/10.4171/owr/2023/38}.
\end{itemize}
\label{end:30005626}
\end{document}
